\section{A complete count instrument along the free motion}\label{sec:instrument}

\subsection{The transported quadrature}
Fix $\gamma>0$, set $T=\gamma^{-1}$, and keep
\begin{equation}\label{eq:Ut}
 U(t)=\ee^{-\ii\omega tL_0},\qquad Q_t=U(t)XU(t)^*.
\end{equation}
The space $\mathscr C=\Sch(\R)\otimes\K_{\rm fin}$, with finite sums
understood, is an invariant core for the required constituent powers.
Conjugating $\ii\partial_k$ by the Fourier multiplier
$\ee^{-\ii\omega t(k^2+4n^2)}$ gives
\begin{equation}\label{eq:Qt}
 Q_t=X-2\omega tP\quad\hbox{on }\mathscr C,\qquad
 \Dom Q_t^q=U(t)\Dom X^q.
\end{equation}
Unitary conjugation specifies the closed realization; no equality with
$\Dom X\cap\Dom P$ is required. Define the strong-law domain
\begin{equation}\label{eq:Dt}
 \mathcal D_t=\Dom L_0\cap\Dom Q_t^2
             =U(t)(\Dom L_0\cap\Dom X^2).
\end{equation}
Both the controller and readout follow $Q_t$. A fixed laboratory-position
attenuation would interact differently with the kinetic dispersion, since
$P^2$ and $X^2$ do not commute.

\subsection{Bounded outcomes, completeness and composition}
For $0\le s\le t\le T$ and $a=\gamma(t-s)$ define
\begin{equation}\label{eq:K}
 B_j(a)=\ee^{-aX^2/2}\frac{(\sqrt aX)^j}{\sqrt{j!}},\qquad
 K_j(t,s)=U(t)B_j(a)U(s)^*,\qquad j=0,1,2,\ldots.
\end{equation}
At $a=0$ put $B_0(0)=I$ and $B_j(0)=0$ for $j>0$. The signed factor
$X^j$ is retained, including for odd counts. The count label is an apparatus
outcome, not an additional angular multiplicity.

\begin{theorem}[Complete ideal instrument]\label{thm:instrument}
Each $K_j(t,s)$ is bounded. For a positive trace-class state $\rho$, the
operations
\[
 \mathcal I_{t,s}(J)(\rho)=\sum_{j\in J}K_j(t,s)\rho K_j(t,s)^*,
 \qquad J\subseteq\mathbb Z_{\ge0},
\]
converge in trace norm and form a countably additive completely positive
instrument. The all-count channel preserves trace. Sequential intervals
aggregate counts exactly, and the all-count channels compose.
\end{theorem}
\begin{proof}
For $a>0$, the scalar squared multiplier is $\ee^{-y}y^j/j!$, $y=av^2$.
Its maximum gives
\begin{equation}\label{eq:Knorm}
 \norm{B_0(a)}=1,\qquad
 \norm{B_j(a)}^2=\frac{j^j\ee^{-j}}{j!}\le1\quad(j\ge1).
\end{equation}
Bounded pointwise convergence proves the strong zero-duration limits.
Operator-norm continuity at zero is neither used nor generally true.
Let $c_N(y)=\ee^{-y}\sum_{j=0}^Ny^j/j!$. The positive exponential series
satisfies $0\le c_N\uparrow1$. For $0\le y\le R_0$ and $N+2\ge2R_0$,
its tail is bounded by $2R_0^{N+1}/(N+1)!$, by a geometric bound on the
successive ratios. On the complementary spectral interval it is at most
one. Consequently
\begin{equation}\label{eq:completeness}
 \sum_{j=0}^NK_j(t,s)^*K_j(t,s)
 =U(s)c_N(aX^2)U(s)^*\xrightarrow{\rm strong}I.
\end{equation}
Spectral decomposition of a positive trace-class $\rho$ and domination by
its summable eigenvalues show that the omitted input trace tends to zero.
Output tails are positive, so their trace norms equal their traces.
This proves trace-norm convergence and trace preservation. The map
$V\psi=(K_j\psi)_{j\ge0}$ is an isometry by \eqref{eq:completeness};
projection onto a subset of count coordinates followed by the count-space
partial trace proves complete positivity and countable additivity, including
finite matrix amplifications.

For $r\le s\le t$, set $a=\gamma(t-s)$, $b=\gamma(s-r)$ and $c=a+b>0$.
Direct multiplication gives, with $J=j+\ell$,
\begin{equation}\label{eq:composition}
 K_j(t,s)K_\ell(s,r)
 =\sqrt{\binom Jj}\left(\frac ac\right)^{j/2}
                     \left(\frac bc\right)^{\ell/2}K_J(t,r).
\end{equation}
The \emph{squared} coefficients sum to one for each $J$. Summing the
corresponding operations, rather than the amplitudes, therefore recovers
the aggregate-count operation. Positive trace-norm convergence justifies
regrouping. At zero durations the stated identity/zero convention applies.
Summing over $J$ gives composition of the all-count channels.
\end{proof}

The instrument is defined by bounded operations on all states. This proof
supplies the required infinite-count result directly; it does not invoke an
unbounded-generator master-equation theorem or a microscopic dilation.

\subsection{The no-count strong law}
\begin{theorem}[No-count propagation]\label{thm:strong}
The no-count operators satisfy $K_0(t,s)K_0(s,r)=K_0(t,r)$, send
$\mathcal D_s$ into $\mathcal D_t$, and obey
\begin{equation}\label{eq:strong}
 \partial_tK_0(t,s)\psi=
 \left[-\ii\omega L_0-\frac\gamma2Q_t^2\right]K_0(t,s)\psi,
 \qquad\psi\in\mathcal D_s.
\end{equation}
For fixed $s<T$, derivatives are right-sided at $t=s$ and left-sided at
$t=T$. The singleton $s=t=T$ has no nontrivial fixed-$s$ derivative
inside the prescribed time interval.
\end{theorem}
\begin{proof}
The cocycle follows from \eqref{eq:composition} at count zero. Put
$\chi=U(s)^*\psi\in\Dom L_0\cap\Dom X^2$ and
$m_a(v)=\ee^{-av^2/2}$. Its first two derivatives have the bounds
\[
 \norm{m_a'}_\infty\le\sqrt{a/\ee},\qquad
 \norm{m_a''}_\infty\le(1+2/\ee)a.
\]
The weak Sobolev product rule gives
\begin{equation}\label{eq:product}
 L_0(m_a\chi)=m_aL_0\chi-2m_a'\partial_v\chi-m_a''\chi,
 \qquad X^2m_a\chi=m_aX^2\chi.
\end{equation}
Since $\chi\in\Dom P^2\cap\Dom A$, all terms belong to $\HH$.
The angular operator commutes with $m_a$. These formulas prove preservation
of both constituent domains and graph continuity, including at $a=0$.
No mixed $XP$ assumption is inserted.
For admissible $a,a+h\ge0$, the mean-value bound
$|(m_{a+h}-m_a)/h|\le v^2/2$ gives
$\partial_a(m_a\chi)=-X^2m_a\chi/2$ strongly on $\Dom X^2$.
The free-unitary derivative holds on $\Dom L_0$ by full spectral calculus.
Combining their difference quotients and using graph continuity proves
\eqref{eq:strong}, together with the conjugated-domain assertion.
The same argument with one-sided quotients handles the endpoints.
\end{proof}

The effective expression is
$H_{\rm kin}-\ii\hbar\gamma Q_t^2/2$, with coupling notation
$\sqrt\gamma Q_t$. Its sign and half factor agree with the standard
no-count convention \cite[Section 1, (4)--(11)]{CastinDalibardMolmer2008v1}.
It is not asserted to be self-adjoint, and the theorem makes no claim about
a larger maximal domain for that formal effective generator. The domain
and convergence statements are exactly what the subsequent preparation
and observation require \cite[Section 4]{TNDDirect}.
